% !TeX root = ../papers/nc-p3.tex
\section{四次曲面の極三次式と十六点での反対座標の一致}
\label{sec:exclude-tetrahedral-quartic}
\begin{lemma}\label{lem:opposite-coincidence-fourteen}
被約二十点の全階数反対平方根 $N$ が一般点で可逆なら、入力と出力の射影線型同型で
対応する印を同定できる個数は高々14である。
\end{lemma}
\begin{proof}
もし15個以上が一致すれば、この同定のもと六次元の交代双一次式はそれらの上で全て消える。
残る高々五点を通る非零な一式 $y^t A x$ を取れる。$N(p)$ は $y(p)y(p)^t$ のスカラー倍だから
$N(x)Ax$ の四つの三次式は $X$ 上で全て消える。三次評価は単射なので恒等零である。
一般点での可逆性から $Ax=0$ が恒等的に成り立ち、矛盾する。
\end{proof}
\begin{theorem}\label{thm:exclude-tetrahedral-quartic}
四辺有限の標準有限Frobenius窓において $X$ は被約で、全階数反対平方根の階数2以下の軌跡が有限なら、
反対八次式には次数5以上の既約因子がある。
\end{theorem}
\begin{proof}
定理\ref{thm:quartic-octic-splitting}の四面体と既約四次式 $F$ の場合を排除すればよい。
$F$ の二重の印を $Z$ とすると $|Z|=16$ である。$Z$ を通る三次式の全空間は四次元であり、
$F$ の核座標の四切断はその基底 $u$ として持ち上がる。実際前定理の $L=P$ を使い、
核像は二重の印二個により非退化で、次数3の多項式に四次の関係はない。
従って定数係数を適切に選んだベクトル $u$ は $F$ の一般点での核を与え、ある一次式ベクトル $b$ により
\[
                            N(x)u(x)=F(x)b(x)
\]
と恒等的に書ける。

$F$ の四偏導は独立である。さもなくば線型座標変更で平面既約四次曲線上の錐となり、
その特異点は頂点と高々三本の母線に限られる。実際平面既約四次の算術種数3は特異点数を上から抑える。
頂点は重複度4ゆえ印でなく、各直線は高々三個の印を含むから16個の二重の印を持てない。
従って $u=B\nabla F$ となる可逆定数行列 $B$ がある。

各 $p\in Z$ で接二次曲線への核写像は閉埋込みで、その像は平面 $y(p)^\perp$ 全体を張る。
ここでの核座標は、四つの三次式の $p$ での一次部、すなわち $B\operatorname{Hess}_p(F)$ によって与えられる。
実際吹上げで例外因子を一回引いた三次式を例外平面に制限している。
従ってHessianは階数3である（Euler式より $p$ はその核に入る）。特に十六点は通常二重点である。
上の恒等式を $p$ で微分すると
\[
                   N(p)B\operatorname{Hess}_p(F)=0 .
\]
入力の同次四空間上でHessianの左核は $p^t$ の張る直線であり、
$N(p)$ は非零の $y(p)y(p)^t$ に比例する。従って全十六点について $[B^t y(p)]=[p]$ となる。
これは補題\ref{lem:opposite-coincidence-fourteen}に反する。
\end{proof}
